\documentclass[12pt,a4paper]{article}

% ── Packages ──────────────────────────────────────────────────────────────────
\usepackage[margin=2.5cm]{geometry}
\usepackage{amsmath, amssymb, amsthm}
\usepackage{enumitem}
\usepackage{tikz}
\usepackage{pgfplots}
\pgfplotsset{compat=1.18}

% ── Macros ─────────────────────────────────────────────────────────────────────
\newcommand{\R}{\mathbb{R}}
\newcommand{\ds}{\displaystyle}

% ── Title ──────────────────────────────────────────────────────────────────────
\title{\textbf{Quiz 1 -- Calculus I (Class 64)}}
\author{}
\date{}

\begin{document}
\maketitle

% ==============================================================================
%  PROBLEMS
% ==============================================================================

\begin{enumerate}

  \item Solve the inequality $3 \le |2x - 5| \le 7$.
        Express the solution as a \emph{set} notation and an \emph{interval} notation.

  \item Find the coordinate $(a, b)$ if a line passes through $(-2, 0)$ and
        $(a, b)$ has gradient $2$, and a line passes through $(3, 5)$ and
        $(a, b)$ has gradient $3/4$.

  \item Given two functions $\ds f(x) = \sqrt{x^2 - 9}$ and $\ds g(x) = \dfrac{1}{x}$.
        \begin{enumerate}
          \item Find $\ds\left(\frac{f}{g}\right)(x)$ and its domain.
          \item Find $(f \circ g)(x)$ and its domain.
        \end{enumerate}

  \item Given $f(x) = 2 + \sqrt{25 - x^2}$ with $0 \le x \le 5$.
        \begin{enumerate}
          \item Find $f^{-1}(x)$ and its domain.
          \item Sketch the graphs of $f(x)$ and $f^{-1}(x)$ in one coordinate system.
        \end{enumerate}

\end{enumerate}

\newpage

% ==============================================================================
%  SOLUTIONS
% ==============================================================================

\begin{enumerate}

  % ------ Problem 1 -----------------------------------------------------------
  \item We split the compound absolute-value inequality $3 \le |2x-5| \le 7$
        into two cases.

        \begin{enumerate}
          \item $2x - 5 \ge 0$ (i.e.\ $x \ge 5/2$):
                \[
                  3 \le 2x - 5 \le 7
                  \;\Longrightarrow\;
                  8 \le 2x \le 12
                  \;\Longrightarrow\;
                  4 \le x \le 6.
                \]

          \item $2x - 5 < 0$ (i.e.\ $x < 5/2$):

                The inequality becomes $3 \le -(2x-5) \le 7$,
                i.e.\ $3 \le 5-2x \le 7$:
                \[
                  3 \le 5 - 2x \le 7
                  \;\Longrightarrow\;
                  -2 \le -2x \le 2
                  \;\Longrightarrow\;
                  -1 \le x \le 1.
                \]
        \end{enumerate}

        Taking the union of both cases:
        \[
          \boxed{S = \{x \in \R \mid -1 \le x \le 1 \text{ or } 4 \le x \le 6\}
            = [-1,\,1]\cup[4,\,6].}
        \]

        % ------ Problem 2 -----------------------------------------------------------
  \item Let $(a, b)$ be the unknown point.

        From the gradient conditions:
        \begin{enumerate}
          \item Line through $(-2, 0)$ and $(a, b)$ has gradient $2$:
                \[
                  \frac{b - 0}{a - (-2)} = 2
                  \;\Longrightarrow\;
                  b = 2a + 4. \tag{1}
                \]

          \item Line through $(3, 5)$ and $(a, b)$ has gradient $3/4$:
                \[
                  \frac{b - 5}{a - 3} = \frac{3}{4}
                  \;\Longrightarrow\;
                  4b = 3a + 11. \tag{2}
                \]
        \end{enumerate}

        Substituting (1) into (2):
        \[
          4(2a + 4) = 3a + 11
          \;\Longrightarrow\;
          8a + 16 = 3a + 11
          \;\Longrightarrow\;
          a = -1,\quad b = 2.
        \]
        \[
          \boxed{(a, b) = (-1,\; 2).}
        \]

        \textit{Verification:}
        \begin{itemize}
          \item Gradient $(-2,0)\to(-1,2)$: $\dfrac{2}{1} = 2$. \checkmark
          \item Gradient $(3,5)\to(-1,2)$: $\dfrac{-3}{-4} = \dfrac{3}{4}$. \checkmark
        \end{itemize}

        % ------ Problem 3 -----------------------------------------------------------
  \item Given $f(x)=\sqrt{x^2-9}$ and $g(x)=\dfrac{1}{x}$.

        Preliminary domains:
        \begin{itemize}
          \item $D_f = (-\infty,-3]\cup[3,+\infty)$ (need $x^2-9\ge0$).
          \item $D_g = \R\setminus\{0\}$ (need $x\ne0$).
        \end{itemize}

        \begin{enumerate}
          \item $\left(\dfrac{f}{g}\right)(x)$:
                \[
                  \left(\frac{f}{g}\right)(x)
                  = \frac{\sqrt{x^2-9}}{\dfrac{1}{x}}
                  = x\sqrt{x^2-9}.
                \]
                Since $x\le-3$ or $x\ge3$ already excludes $x=0$:
                \[
                  \boxed{\left(\frac{f}{g}\right)(x) = x\sqrt{x^2 - 9},
                  \qquad D_{f/g} = (-\infty,-3]\cup[3,+\infty).}
                \]

          \item $(f \circ g)(x)$:
                \[
                  (f \circ g)(x) = f\!\left(\frac{1}{x}\right)
                  = \sqrt{\frac{1}{x^2} - 9}
                  = \frac{\sqrt{1 - 9x^2}}{|x|}.
                \]
                Domain requires $x\ne0$ and $g(x)\in D_f$,
                i.e.\ $\dfrac{1}{x}\le-3$ or $\dfrac{1}{x}\ge3$:
                \begin{itemize}
                  \item $\dfrac{1}{x}\ge3$: $\dfrac{1-3x}{x}\ge0
                          \Rightarrow 0<x\le\dfrac{1}{3}$.
                  \item $\dfrac{1}{x}\le-3$: $\dfrac{1+3x}{x}\le0
                          \Rightarrow -\dfrac{1}{3}\le x<0$.
                \end{itemize}
                \[
                  \boxed{(f\circ g)(x) = \frac{\sqrt{1-9x^2}}{|x|},
                    \qquad D_{f\circ g} =
                    \left[-\tfrac{1}{3},\,0\right)\cup\left(0,\,\tfrac{1}{3}\right].}
                \]
        \end{enumerate}

        % ------ Problem 4 -----------------------------------------------------------
  \item Given $f(x) = 2 + \sqrt{25 - x^2}$, $0 \le x \le 5$.

        \begin{enumerate}
          \item Let $y = 2 + \sqrt{25 - x^2}$. Solving for $x$:
                \begin{align*}
                  y - 2   & = \sqrt{25 - x^2},    \\
                  (y-2)^2 & = 25 - x^2,           \\
                  x^2     & = 25 - (y-2)^2,       \\
                  x       & = \sqrt{25 - (y-2)^2}
                  \quad (\text{positive root since } x \ge 0).
                \end{align*}
                Swapping $x \leftrightarrow y$:
                \[
                  \boxed{f^{-1}(x) = \sqrt{25 - (x-2)^2}.}
                \]
                Domain of $f^{-1}$ $=$ Range of $f$. Since $f(0)=7$, $f(5)=2$,
                and $f$ is strictly decreasing: $D_{f^{-1}} = [2,\,7]$.

          \item Both curves are quarter-circle arcs (radius $5$) symmetric
                about the line $y = x$.

                \begin{center}
                  \begin{tikzpicture}[scale=1.1]
                    % Axes
                    \draw[->] (-0.3,0) -- (7.8,0) node[right] {$x$};
                    \draw[->] (0,-0.3) -- (0,7.8) node[above] {$y$};

                    % Tick marks and labels
                    \foreach \x in {1,...,7} {
                        \draw (\x,2pt) -- (\x,-2pt) node[below] {\small $\x$};
                      }
                    \foreach \y in {1,...,7} {
                        \draw (2pt,\y) -- (-2pt,\y) node[left] {\small $\y$};
                      }

                    % y = x (dashed)
                    \draw[dashed, gray] (0,0) -- (7.5,7.5);

                    % f(x) = 2 + sqrt(25 - x^2), x in [0,5]
                    \draw[blue, thick, domain=0:5, smooth, samples=120]
                    plot (\x, {2 + sqrt(25 - \x*\x)});
                    \node[blue, above right] at (3,6) {$f(x)$};

                    % f^{-1}(x) = sqrt(25 - (x-2)^2), x in [2,7]
                    \draw[red, thick, domain=2:7, smooth, samples=120]
                    plot (\x, {sqrt(25 - (\x-2)*(\x-2))});
                    \node[red, above right] at (6,3) {$f^{-1}(x)$};

                    % Key points — f(x)
                    \filldraw[blue] (0,7) circle (1.8pt);
                    \filldraw[blue] (5,2) circle (1.8pt);
                    \node[blue, below right] at (5,2)
                    {\small $(5,2)$};

                    % Key points — f^{-1}(x)
                    \filldraw[red] (2,5) circle (1.8pt);
                    \node[red, below left] at (2,5)
                    {\small $(2,5)$};
                    \filldraw[red] (7,0) circle (1.8pt);

                    % y=x label
                    \node[gray, above right] at (5.5,5.5)
                    {\small $y=x$};
                  \end{tikzpicture}
                \end{center}
        \end{enumerate}

\end{enumerate}

\end{document}
